\documentclass[12pt]{article}
\usepackage[utf8]{inputenc}
\usepackage[T1]{fontenc}
\usepackage{hyperref}
\title{Berawan, West Language Resources for K-12 Curriculum Development}
\author{Compiled by Claude Sonnet 5.5 (claude-sonnet-5-5) for Andr\'as Kornai}
\date{}
\begin{document}
\maketitle

\section{Introduction}
West Berawan (ISO 639-3 zbw) is an Austronesian language of the North Sarawakan group, spoken in eastern Sarawak, Malaysia. Ethnologue gives about 1,300 speakers, with Long Terawan village near Mulu as its location. The prompt gives about 1,000. Joshua Project uses the name Long Pata for the group. Sources do not agree on whether ``West Berawan'' is the same as the Long Terawan dialect or a broader grouping. A 2024 source calls all the Berawan lects endangered. Little material exists and none for schools was found.

\section{Primary Educational Resources}
None found. The language is not known to be taught in schools.

\section{Linguistic and Reference Materials}
\begin{itemize}
\item Wikipedia article on the Berawan language: \url{https://en.wikipedia.org/wiki/Berawan_language}
\item Austronesian Comparative Dictionary language page: \url{https://acd.clld.org/languages/19364}
\item PARADISEC catalogue item (Blust recordings collection, RB1): \url{https://catalog.paradisec.org.au/repository/RB1/003}
\item Article on Berawan lects including Long Terawan, Borneo Research Journal: \url{https://publisher.unimas.my/ojs/index.php/BJK/article/view/2411/1012}
\item Further UNIMAS journal article from the same search: \url{https://publisher.unimas.my/ojs/index.php/BJK/article/view/8818/2484}
\end{itemize}
Clayre and Wan (1996) on the phonology of the Long Terawan dialect, and Blust (1984), were mentioned in search summaries, but I did not see a URL for either.

\section{Religious and Community Resources}
\begin{itemize}
\item Joshua Project profile of the Western Berawan in Malaysia: \url{https://joshuaproject.net/people_groups/10800}
\item Sarawak Museum page (content not examined): \url{https://museum.sarawak.gov.my/web/subpage/webpage_view/611}
\item Sarawak Museum page (content not examined): \url{https://museum.sarawak.gov.my/web/subpage/webpage_view/847}
\end{itemize}

\section{Digital and Audio Resources}
The PARADISEC item above is the only audio-related record found.

\section{Curriculum Development Context}
West Berawan is spoken by all adults but not all young people, and it is not taught in schools. A curriculum would start with documentation: the existing phonology and wordlists, and new recordings made with community members. Malay and English would remain the languages of instruction, and Berawan could be added as a local subject. A lack of agreement on what counts as West Berawan would need settling with the community first.

\end{document}